\documentclass[11pt]{article}

% ---------- Layout ----------
\usepackage[margin=0.8in, top=0.6in, bottom=0.7in, headsep=10pt]{geometry}
\usepackage{times}
\usepackage[skip=3pt]{parskip}
\usepackage{titlesec}
\usepackage{enumitem}
\usepackage{fancyhdr}
\usepackage{hyperref}
\usepackage{graphicx}
\usepackage{amsmath}
\usepackage{booktabs}
\usepackage{float}
\usepackage{xcolor}

\hypersetup{
    colorlinks=false,
    pdfborder={0 0 0},
    pdftitle={Week 3 Progress Report},
    pdfauthor={Aryan Thakkar, Nikkesh Parekh, Hitaksh Thaker, Maitri Parekh}
}

% ---------- Section styling ----------
\titleformat{\section}
  {\normalfont\large\bfseries}
  {\thesection}{0.6em}{}
  
\titleformat{\subsection}
  {\normalfont\normalsize\bfseries}
  {\thesubsection}{0.5em}{}

\titlespacing*{\section}{0pt}{6.5pt}{2.5pt}
\titlespacing*{\subsection}{0pt}{5pt}{1.5pt}

\renewcommand{\thesection}{\arabic{section}}

% ---------- List spacing ----------
\setlist[itemize]{
    itemsep=1pt,
    topsep=2pt,
    parsep=0pt,
    leftmargin=1.4em
}

\setlist[enumerate]{
    itemsep=1pt,
    topsep=2pt,
    parsep=0pt,
    leftmargin=1.6em
}

% ---------- Header / footer ----------
\pagestyle{fancy}
\fancyhf{}

\renewcommand{\headrulewidth}{0.4pt}
\renewcommand{\footrulewidth}{0.4pt}

\fancyhead[L]{\small\textsc{ECE 501 --- Digital Image Processing}}
\fancyhead[R]{\small Week 3 Progress Report}

\fancyfoot[L]{\small\textit{Galaxy Morphology Classification --- Week 3}}
\fancyfoot[R]{\small\thepage}

% ============================================================
% DOCUMENT
% ============================================================

\begin{document}

\thispagestyle{fancy}

% ============================================================
% TITLE BLOCK
% ============================================================

\begin{center}

\Large\textbf{Morphological Classification of Galaxies in HST and JWST Images Using Classical DIP}

\vspace{3pt}

\normalsize\textbf{Weekly Progress Report --- Week 3}

\vspace{4pt}

\small
\textbf{Aryan Thakkar}
\quad $\cdot$ \quad
\textbf{Nikkesh Parekh}
\quad $\cdot$ \quad
\textbf{Hitaksh Thaker}
\quad $\cdot$ \quad
\textbf{Maitri Parekh}

\vspace{3pt}

\small ECE 501 -- Digital Image Processing Course Project

\vspace{2pt}

\small\textbf{Faculty:} Prof. Mehul Raval
\quad $\cdot$ \quad
\textbf{TAs:} Apoorva Mehta \quad $\cdot$ \quad Jyoti Triklani \quad $\cdot$ \quad Sagar Chavda

\vspace{3pt}

\small Date: \today

\end{center}

\vspace{1pt}
\hrule height 0.8pt
\vspace{3pt}

% ============================================================
% 1. INTRODUCTION AND OBJECTIVES
% ============================================================

\section{Week 3 Scope}
Week 2 established the literature baseline, defined the EDA parameters, identified the dataset sources, and outlined the algorithmic pipeline direction. Week 3 executes the first concrete technical phase of that plan: a rigorous, comparative R\&D study of classical restoration and enhancement methods applied to the acquired JWST NIRCam and HST ACS/WFC FITS data for Stephan's Quintet (HCG 92). Morphological segmentation and classification are deliberately deferred to Week 4; this report covers only image restoration, denoising, contrast enhancement, compositing, and quantitative dataset evaluation.

\section{Deconvolution Experiments}
PSF-induced optical blur was identified in Week 2 as a primary degradation source. Two deconvolution strategies were implemented on a $500 \times 500$ crop of the JWST NIRCam F200W channel centred on the NGC 7318A/B interacting pair, using an approximated $7 \times 7$ Gaussian PSF ($\sigma = 1.5$). Sharpness was quantified using the variance of the image Laplacian ($\sigma^2_{\nabla^2}$).

\subsection{Richardson-Lucy (RL) Deconvolution}
An iterative expectation-maximisation algorithm suited to Poisson-distributed photon noise, implemented via \texttt{skimage.restoration.richardson\_lucy}. Evaluated at 5, 20, and 50 iterations. Laplacian variance increased monotonically with iterations, confirming progressive edge sharpening. At 50 iterations, however, background noise amplification became severe and ringing artefacts appeared around bright stellar cores, making high-iteration RL unsuitable for the downstream segmentation phase.

\subsection{Unsupervised Wiener Deconvolution}
Implemented via \texttt{skimage.restoration.unsupervised\_wiener}. Minimises mean-square error between estimated and true image by simultaneously solving the inverse filter and noise-suppression problem. Does not iterate, so noise amplification is bounded. Produced cleaner background relative to RL-20 while preserving galactic-bulge edge definition.

\subsection{Selection}
Wiener deconvolution selected as the pipeline's deconvolution step. It provides a superior noise-sharpness trade-off compared to RL at equivalent visual sharpness, without the ringing artefacts that would corrupt morphological boundary detection in Week 4. The PSF will be refined using the WebbPSF tool in a future iteration.

\section{Denoising Method Comparison}
Five classical denoising filters were applied to the same preprocessed F200W crop and evaluated on two criteria: (1) background noise variance measured in a flat $100 \times 100$ sky patch (lower is better), and (2) qualitative preservation of faint tidal-tail structure.

Bilateral filtering selected (\texttt{cv2.bilateralFilter}, d=9, sigmaColor=75, sigmaSpace=75). It is the only method that simultaneously flattens background noise variance and preserves the sharp intensity transitions at galaxy boundaries required for accurate segmentation thresholding in Week 4.

\section{Contrast Enhancement Comparison}
Astronomical images have extreme dynamic range --- the AGN core of NGC 7319 is orders of magnitude brighter than surrounding tidal tails. Week 2's EDA confirmed this via dynamic range analysis. Five stretch methods were evaluated on the same normalised image to determine which best reveals faint structure without saturating bright cores.

\begin{itemize}
    \item \textbf{Linear Stretch} --- clips faint tidal structures entirely; bright core dominates. Rejected.
    \item \textbf{Logarithmic Stretch} --- faint background galaxies visible; mild mid-tone washout in arms.
    \item \textbf{Global Histogram Equalisation} --- severe core overexposure; amplifies background noise. Rejected.
    \item \textbf{CLAHE} (clip = 0.03) --- strong local contrast but heavy artificial texture in sky background. Rejected.
    \item \textbf{Asinh Stretch} (a = 0.05) --- linear near zero (sky background), logarithmic for bright values (cores). Best dual-range fidelity.
\end{itemize}

Asinh Stretch selected (\texttt{astropy.visualization.AsinhStretch(a=0.05)}). Designed specifically for astronomical data, it preserves both the extreme galactic nuclei brightness and the faint extended interacting arms of NGC 7318A/B simultaneously --- a prerequisite for reliable morphological feature extraction.

\section{Multi-Filter RGB Composite}
JWST NIRCam data was acquired for three filters: F090W (0.9 $\mu$m), F150W (1.5 $\mu$m), and F200W (2.0 $\mu$m). Two filter-to-channel mapping strategies were tested, both using Asinh stretch (a = 0.05) with per-channel 99.5th-percentile clipping:

\begin{itemize}
    \item \textbf{Physical mapping} (R=F200W, G=F150W, B=F090W) --- long-to-short wavelength maps to red-to-blue, consistent with STScI NIRCam standard practice.
    \item \textbf{Inverse mapping} (R=F090W, G=F150W, B=F200W) --- produces blue-dominant image; cool stellar populations appear warm-coloured, increasing inter-galaxy contrast in the interacting region.
\end{itemize}

Physical mapping retained as the primary composite. It is the standard astronomical convention and produces morphologically interpretable colour distributions (warm cores, blue star-forming regions in spiral arms), which supports visual and algorithmic classification in Week 4.

\section{HST vs. JWST Image Quality Evaluation}
To justify the primary dataset selection made in Week 2 as a methodological plan, a quantitative image quality comparison was executed on matched $400 \times 400$ px spatial crops of Stephan's Quintet from both instruments. SNR was defined as image mean divided by the standard deviation of a flat background patch; sharpness as the Laplacian variance; dynamic range as (max - min) of the calibrated pixel range.

Pixel-scale values are published instrument specifications. SNR, sharpness and dynamic range values are computed from the preprocessed FITS crops; exact outputs will be confirmed upon final notebook execution. JWST's 6.5 m primary mirror and near-infrared sensitivity yield substantially higher SNR and finer spatial resolution. 

JWST NIRCam selected as the primary dataset for all downstream segmentation and classification work.

\section{Artifact Rejection and Parameter Sensitivity}
\subsection{Cosmic Ray Rejection}
Cosmic ray hits appear as isolated, abnormally bright pixels that would be mistakenly detected as point sources in the segmentation phase. Rejection applied a $5 \times 5$ median filter to produce a smooth reference image, then flagged pixels where the absolute difference exceeded $3\sigma$ of the difference map:

\texttt{artifact\_mask = |I - median\_5x5(I)| > 3 $\times$ $\sigma$\_diff}

Flagged pixels were replaced with the local median value. The $3\sigma$ threshold is conservative --- genuine faint structure is not masked, only sharp isolated spikes.

\subsection{NLM Filter Strength (h) Sweep}
h parameter swept at 3, 10, and 15 to calibrate the smoothing--detail trade-off:
\begin{itemize}
    \item \textbf{h = 3} --- under-smoothed; background noise remains visible in flat regions.
    \item \textbf{h = 10} --- balanced; noise suppressed without blurring spiral arm boundaries.
    \item \textbf{h = 15} --- over-smoothed; fine structural detail in galactic arms lost.
\end{itemize}
Selected: h = 10, templateWindowSize = 7, searchWindowSize = 21.

\subsection{Background2D Box Size Sweep}
\texttt{photutils Background2D} box size swept at 25, 50, and 100 pixels. Box 25 over-fitted, partially subtracting faint galaxy halos. Box 100 under-fitted the non-uniform illumination gradient near the NGC 7320c foreground quasar. Box 50 selected as optimal.

\section{Final Pipeline Configuration}
All method and parameter choices from Sections 2--7 are consolidated into the following pipeline, which will be applied consistently in the Week 4 segmentation and classification phase.

\section{Week 3 Summary and Status}
The Week 3 R\&D phase has successfully formulated a mathematically justified, optimized image processing pipeline. With the foundational images now optimally restored and enhanced, the project is fully prepared to execute classical morphological classification in the subsequent phase.

\section{Planned Work --- Week 4}
\begin{itemize}
    \item Execute the full notebook (Restart \& Run All) and populate Table 1 with actual measured SNR, sharpness, and dynamic range values.
    \item \textbf{B6:} Galaxy segmentation --- adaptive thresholding, connected-component labelling, and star--galaxy separation using the concentration-index criterion developed in the Week 3 pipeline.
    \item \textbf{B7:} CAS feature extraction and rule-based morphological classification (Elliptical / Spiral / Irregular) per detected galaxy.
    \item \textbf{C1:} Repeat segmentation and classification on the HST dataset for the same target and compare classification accuracy against JWST results.
    \item \textbf{B8/C3:} Validate against Galaxy Zoo reference labels for the five primary Stephan's Quintet members (NGC 7317, NGC 7318A, NGC 7318B, NGC 7319, NGC 7320).
\end{itemize}

\end{document}
